\section{三つの見方を結ぶ}\label{節：三つの見方を結ぶ}
\subsection{一つの例を三通りに記録する}

特に
\[
 M=\begin{pmatrix}1&1\\-1&1\end{pmatrix}
\]
を使おう.
状態の更新は$x_{n+1}=x_n+y_n,\,y_{n+1}=-x_n+y_n$である.
複素数を$z_n=y_n+ix_n$と選ぶと,
\[
 z_{n+1}=(1+i)z_n.
\]
この座標では,\,毎回$\sqrt2$倍されながら反時計回りに$\frac{\pi}{4}$回転する.
各成分は特性方程式$X^2-2X+2=0$に対応し,
\[
 x_{n+2}=2x_{n+1}-2x_n
\]
などを満たす.

一方,\,成分比$a_n=\frac{x_n}{y_n}$を記録すると,
\[
 a_{n+1}=\frac{a_n+1}{1-a_n}.
\]
共通の拡大率が消え,\,傾きの変化だけが残る.
さらに$w_n=\Phi(a_n)$で円周へ移すと,
\[
 w_{n+1}=iw_n.
\]
角度を2倍して記録する$\Phi$によって,\,回転角は$\frac{\pi}{2}$になる.

$(x_1,y_1)=(\frac12,1)$なら,\,成分比は先ほどの4周期の列となる.
行列では4回後の状態は$M^4\boldsymbol v_1=-4\boldsymbol v_1$であり,\,元の状態そのものには戻らない.
しかし両成分が同じ倍率を受けるため比は戻る.
\textbf{状態の周期と,\,比の周期は同じとは限らない}.


\begin{bookpoint}
同じ行列から,連立漸化式,隣接3項間漸化式,一次分数型を作れる.
状態全体を追う場合と,成分の一つや比だけを追う場合とでは,読み取れる性質が異なる.
\end{bookpoint}
